\documentclass[aspectratio=169,notheorems]{corvinusmetropolis}
\newtheorem{theorem}{{Tétel.}}
\newtheorem{proposition}[theorem]{{Állítás.}}
\newtheorem{lemma}[theorem]{{Lemma.}}
\newtheorem{FA}[theorem]{{Az algebra alaptétele.}}
\newtheorem{corollary}[theorem]{{Következmény.}}
\theoremstyle{definition}
\newtheorem{definition}[theorem]{{Definíció.}}
\newtheorem{defprop}[theorem]{{Definíció-Állítás.}}
\newtheorem{notes}[theorem]{{Megjegyzés.}}

\email{magyarkuti@uni-corvinus.hu}
\author[Magyarkuti]{Gyula~Magyarkuti\,\orcidlink{0000-0003-4444-0481}}
\institute[Corvinus Egyetem]%
{%\inst{1}
  Budapesti Corvinus Egyetem\\
  Matematika Tanszék\\
  Sóház, 208/B

  \begin{minipage}[b]{0.8\linewidth}
  email: \putemail\\
  letölthető: \url{https://magyarkuti.github.io/algebra-1}
  \end{minipage}
  \hfill{\adjustbox{valign=B}{\qrcode[hyperlink]{https://magyarkuti.github.io/algebra-1}}}%% iff qrcode needed, otherwise comment this line!
}
\date[Tavasz]{\ontoday}%% overwrite later, if needed.

\titlegraphic{\hfill{\includegraphics[width=.2\linewidth]{\myinstlogo}}}

\renewcommand{\emph}{\alert}
\renewcommand{\mathbf}{\mathbb}
\newcommand{\lin}{\operatorname{lin}}
\newcommand{\gen}{\operatorname{gen}}
\newcommand{\im}{\operatorname{Im}}
\newcommand{\id}{\operatorname{id}}
\newcommand{\sgn}{\operatorname{sgn}}
\newcommand{\tr}{\operatorname{tr}}
\newcommand{\codim}{\operatorname{codim}}
%newcommand{\deg}{\operatorname{deg}}
\newcommand{\ip}[2]{\langle#1,#2\rangle}
\newcommand{\uj}{\text{új}}
\newcommand{\rgi}{\text{régi}}
\newcommand{\Star}[1]{#1\ensuremath{^*}\kern-\scriptspace}
\newcommand{\CStar}{\Star{\ensuremath{\mathrm{C}}}}
\endinput

\title[Algebra 1.]{Speciális polinom faktorizáció}
\subtitle{Az algebra alaptételének következményei}
\date[6. hét]{2026.~október 18-án}
\begin{document}
\openingframes{Tartalom}
\section{Polinom faktorizáció a komplex- és a valós számtest felett}
\frame{
A komplex számtest felett minden nem konstans polinomnak van gyöke.
Formálisabban:
\metroset{block=fill}
\begin{FA}
	Minden legalább elsőfokú komplex együtthatós polinomnak van komplex gyöke.
	Azaz
	ha $p\left( t \right)\in\mathbb{C}[t]$ egy nem konstans polinom,
	akkor létezik $z\in\mathbb{C}$ komplex szám,
	amelyre $p\left( z \right)=0$.
\end{FA}
\metroset{block=transparent}
Az állítást itt nem tudjuk igazolni és egyelőre érdemes bizonyítás nélkül elfogadni.
Illusztrációként megértettük, hogy miért igaz $p\left( t \right)=\alpha_0+t^n$ alakú polinomra.
A felépítés jelen pontján tekintsük a komplex számtest egy még nem igazolt tulajdonságának.
}
\frame{
Most összefoglaljuk, hogy az algebra alaptétele mit jelent a komplex test feletti, majd a valós test feletti polinomok
faktorizációjára nézve.
Láttuk, hogy minden normált polinom előáll mint normált irreducibilis polinomok szorzata,
emiatt azt kell meggondolnunk, hogy mik az irreducibilis polinomok.
\metroset{block=fill}
\begin{proposition}
	A $\mathbb{C}\left[ t \right]$ polinomgyűrűben minden legalább másodfokú polinom reducibilis.
\end{proposition}
\metroset{block=transparent}
}
%\begin{proof}
%	Legyen $p$ egy legalább másodfokú komplex polinom.
%	Az algebra alaptétele miatt van gyöke, és
%	tudjuk hogy a gyöktényező mindig kiemelhető.
%	Így azt kapjuk, hogy $p(t)=\left( t-z \right)h\left( t \right)$ alakú,
%	emiatt $2\leq\deg p=1+\deg h$, ergo
%	$p$ valóban előállt mint két legalább elsőfokú polinom szorzata.
%\end{proof}
\subsection{Komplex eset}
\frame{
A polinomok szorzattá bontásáról szóló tételnek a komplex számtest feletti speciális esete tehát:
\metroset{block=fill}
\begin{proposition}\label{pr:PolFact}
	A $\mathbb{C}[t]$ polinomgyűrű minden legalább elsőfokú normált polinomjához léteznek
	a sorrendjüktől eltekintve egyetlen $z_1,\dots,z_s\in\mathbb{C}$ egymástól különböző komplex számok,
	és léteznek $n_1,\dots,n_s$ pozitív egészek,
	amelyekre
	\[
		p\left( t \right)=
		\left( t-z_1 \right)^{n_1}
		\cdots
		\left( t-z_s\right)^{n_s}.\qedhere
	\]
\end{proposition}
\metroset{block=transparent}
Persze ez azt jelenti, hogy egy $n$-ed fokú komplex polinomnak mindig pontosan $n$ komplex gyöke van,
ha a gyökök számát multiplicitással számoljuk.
}
\frame{
Az állítást egy kicsit egyszerűbben úgy is fogalmazhatjuk,
hogy
\emph{
	minden nem konstans komplex polinom előáll mint első fokú komplex polinomok szorzata.
}

Az iménti mondat nyilván nem igaz a ,,komplex'' szót a ,,valós'' szóra cserélve,
ugyanis minden negatív diszkriminánsú másodfokú valós polinom jó is ellenpéldának.
Ez a jelenség az oka annak,
hogy a komplex számtest felett sokkal kényelmesebb dolgoznunk mint a valós számok felett.
}
\subsection{Valós eset}
\frame{
Most nézzük, hogy mit jelent az algebra alaptétele a valós test feletti polinomgyűrűre nézve.
\begin{proposition}
	Legyen $p\left( t \right)\in\mathbb{C}[t]$ a komplex polinomgyűrű egy olyan eleme,
	amelynek minden együtthatója valós.
	Ekkor a $z\in \mathbb{C}$ komplex szám pontosan akkor gyöke $p$-nek,
	ha $\bar{z}$ is gyöke $p$-nek.
	Sőt, ha $z$ egy $k$-szoros multiplicitású gyök,
	akkor $\bar{z}$ is pontosan $k$-szoros multiplicitású gyök.
\end{proposition}
}
%\begin{proof}
%	A $p$ polinomra tehát
%	\begin{math}
%		p\left( t \right)
%		=
%		\sum_{j=0}^n\alpha_jt^j,
%	\end{math}
%	ahol $\alpha_0,\alpha_1,\dots,\alpha_n\in\mathbb{R}$.
%	A konjugálás tulajdonságai szerint egy $z$ komplex számra
%	\[
%		p\left( z \right)=
%		\overline{
%			\sum_{j=0}^n\alpha_jz^j
%		}
%		=
%		\sum_{j=0}^n\overline{\alpha_jz^j}
%		=
%		\sum_{j=0}^n\bar{\alpha_j}\bar{z}^j
%		=
%		\sum_{j=0}^n\alpha_j\bar{z}^j
%		=
%		p\left(\bar{z} \right).
%	\]
%	Emiatt valós együtthatós $p$ polinomra és $z$ komplex számra
%	$p\left( z \right)=0$ akkor és csak akkor, ha $p\left( \bar{z} \right)=0.$
%
%	Ha a gyök speciálisan valós szám, akkor a multiplicitásra vonatkozó állítás semmit mondó.
%	Ha most $z$ egy nem valós komplex gyök,
%	akkor $p(t)=\left( t-z \right)\left( t-\bar{z} \right)h\left( t \right)$, ahol
%	$h\left( z \right)=0$ pontosan akkor, ha $h\left( \bar{z} \right)=0$.
%	Így, ha $z$ egy $k$ szoros gyök,
%	akkor
%	\[
%		p\left( t \right)=\left( t-z \right)^{k}\left( t-\bar{z} \right)^{k}\cdot h\left( t \right)
%	\]
%	alakú, ahol $h$-nak már sem $z$ sem $\bar{z}$ nem gyöke.
%\end{proof}
\metroset{block=fill}
\frame{
Minden valós együtthatós polinom előáll mint első vagy másodfokú polinomok szorzata:
\begin{proposition}\label{pr:RealPolFact}
	Az $\mathbb{R}[t]$ polinomgyűrű minden legalább elsőfokú normált polinomjához
	léteznek $x_1,\dots,x_r\in\mathbb{R}$ valós számok,
	léteznek $\alpha_1,\beta_1,\dots,\alpha_s,\beta_s$ valós együttható párok,
	és léteznek
	$n_1,\dots,n_r, m_1,\dots,m_s$ pozitív egészek úgy, hogy
	\[
		p\left( t \right)
		=
		\left( t-x_1 \right)^{n_1}
		\cdots
		\left( t-x_r \right)^{n_r}
		\cdot
		\left(\alpha_1 + \beta_1t +t^2\right)^{m_1}
		\cdots
		\left(\alpha_s + \beta_st +t^2\right)^{m_s}.
	\]
	Itt $r,s\geq 0$, de $r+s>0$, és $n_1+\dots+n_r+2\left( m_1+\dots+m_s \right)=\deg p$,
	továbbá a jobboldalon szereplő másodfokú polinomok irreducibilisek.
\end{proposition}
}
\metroset{block=transparent}
%\begin{proof}
%	Tekintsük a $p$ polinomot mint a komplex számtest feletti polinomgyűrű egy elemét,
%	és alkalmazzuk a komplex polinom faktorizációról szóló tételt.
%	A $p$ így előáll mint első fokú, esetleg komplex polinomok szorzata.
%	A gyököket osszuk két része.
%	Legyenek $x_1,\dots,x_r$ a különböző valós gyökök, amelyek multiplicitása rendre
%	$n_1,\dots,n_r$.
%	Legyenek $z_1,\overline{z_1},\dots,z_s,\overline{z_s}$ a különböző nem valós de komplex gyökök,
%	ahol $m_1,\dots,m_s$ rendre a konjugált gyökpárok multiplicitása.
%	Így
%	\[
%		p\left( t \right)=
%		\left( t-x_1 \right)^{n_1}
%		\cdots
%		\left( t-x_r \right)^{n_r}
%		\cdot
%		\left( \left( t-z_1 \right)\left( t-\bar{z_1} \right) \right)^{m_1}
%		\cdots
%		\left( \left( t-z_s \right)\left( t-\bar{z_s} \right) \right)^{m_s}.
%	\]
%	Világos, hogy $r,s\geq 0$, de $r+s>0$, és $n_1+\dots+n_r+2\left( m_1+\dots+m_s \right)=\deg p$.
%	A komplex faktorokra végezzük el a szorzást,
%	így
%	$
%		\left( t-z_k \right)\left( t-\bar{z_k} \right)
%		=
%		t^2-2\Re z_kt+|z_k|^2
%	$.
%	Így $\alpha_k=|z_k|^2$ és $\beta_k=-2\Re z_k$, választással készen is vagyunk.
%\end{proof}
\frame{
A fenti állításból két dolog azonnal látszik.
Az első,
hogy \alert{\emph{valós számtest felett minden legalább harmadfokú polinom reducibilis}},
a második pedig,
hogy ha egy valós együtthatós polinomnak nincs valós gyöke,
akkor csak páros fokú lehet, vagy ami ugyanaz:
\alert{\emph{minden páratlan fokú valós együtthatós polinomnak van valós gyöke.}}
}

\closingframes{Köszönöm a figyelmet, a türelmet, a megértést, az empátiát!}
\end{document}
